\section{From uniform aisle inclusions to a slicing}
\label{sec:grid}

We now use the bounded aisles $U_\eta$ and their hearts
$\mathcal A_\eta$ constructed above.  The parameter $\eta=(\eta_0,\eta_1)$
ranges over the fixed open square $\Omega$ of
Proposition~\ref{prop:uniform-aisle}, and $0\in\Omega$.  The inputs to this
section are
\begin{equation}\label{eq:grid-inputs}
 \begin{gathered}
 \Phi U_\eta\subset U_\theta
 \quad(\eta,\theta\in\Omega),\\
 N_\eta=d+\tfrac12c+ur+\eta_1c+\eta_0r\ge0
 \quad\text{on }\mathcal A_\eta,
 \end{gathered}
\end{equation}
the absence of nonzero objects of zero numerical class in these hearts,
and Lemma~\ref{lem:periodicity}, Lemma~\ref{lem:numerical}, and
Proposition~\ref{prop:central-charge}.  In particular, throughout this
section $\Phi$ is an autoequivalence, $\Phi^5\simeq[2]$, and
$Z\Phi=e^{2\pi i/5}Z$.  All extension closures below mean finite extension
closures, contain zero, and are strictly full subcategories of $\D$.
For a subcategory $\mathcal C$, write
$\mathcal C^\perp=\{E:\Hom(C,E)=0\text{ for all }C\in\mathcal C\}$;
the left orthogonal ${}^\perp\mathcal C$ is defined similarly.

The decisive use of the perturbations is twofold.  Simultaneous positivity
in fixed middle hearts improves the coarse nesting to a grid with spacing
$1/5$.  Positivity in every perturbed preceding heart then bounds the full
numerical class of each grid block by its charge.

\subsection{The finer grid}

Put $D_i^\eta=\Phi^iU_\eta$ for $i\in\Z$.  The hypotheses give
\begin{equation}\label{eq:coarse-grid}
 D_{i+1}^\eta\subset D_i^\theta\quad
 (i\in\Z,\ \eta,\theta\in\Omega),\qquad
 D_{i+5}^\eta=D_i^\eta[2].
\end{equation}
All successive inclusions allow independent parameter choices.
The charge rotation suggests how to insert a fifth-step between these
coarse cuts: $\Phi^3[-1]$ multiplies $Z$ by $e^{i\pi/5}$.
For its translates to give nested aisles, we need
$\Phi^3U_\alpha[-1]\subset U_\beta$, or equivalently
$D_3^\alpha\subset D_0^\beta[1]$.  The next proposition proves this
inclusion before any slicing or semistable phase is defined.

\begin{proposition}[Uniform refinement of the grid]\label{prop:fine-grid}
There is an open square $\Omega_1$ about zero, contained in $\Omega$, such
that
\begin{equation}\label{eq:single-shift-inclusion}
 D_3^\alpha\subset D_0^\beta[1]
 \qquad(\alpha,\beta\in\Omega_1).
\end{equation}
For $k\in\Z$ and $\eta\in\Omega_1$, define
\begin{equation}\label{eq:fifths-aisles}
 V_k^\eta=\Phi^mU_\eta[j],\qquad k=2m+5j.
\end{equation}
This is independent of the integers $m,j$ chosen, and these bounded aisles
satisfy
\begin{equation}\label{eq:fifths-nesting}
 \begin{gathered}
 V_{k+1}^\alpha\subset V_k^\beta,\qquad
 V_{k+5}^\eta=V_k^\eta[1],\qquad
 \Phi V_k^\eta=V_{k+2}^\eta,\\
 k\in\Z,\quad \alpha,\beta,\eta\in\Omega_1.
 \end{gathered}
\end{equation}
\end{proposition}

\begin{proof}
For $\alpha,\beta\in\Omega$, consider
\[
 \mathcal I_{\alpha,\beta}
   =D_3^\alpha\cap(D_0^\beta[1])^\perp.
\]
Every object of this intersection belongs to the hearts of $D_0^\beta$
and $D_3^\alpha$.  More strongly, for every $\gamma\in\Omega$ and
$i\in\{1,2\}$, \eqref{eq:coarse-grid} gives
\[
 D_3^\alpha\subset D_i^\gamma\subset D_0^\beta,
 \qquad D_i^\gamma[1]\subset D_0^\beta[1].
\]
The first inclusion gives aisle membership, and the second gives the
orthogonal membership.  Thus the same object lies in the heart of
$D_i^\gamma$ for every $\gamma$ in the original fixed square $\Omega$.
This square is independent of the endpoints $\alpha,\beta$ and of the
object.

Fix a Euclidean norm on $\Knum(X)_\R$.  Suppose there were nonzero
objects $E_n\in\mathcal I_{\alpha_n,\beta_n}$ with
$\alpha_n,\beta_n\longrightarrow0$.  The endpoint heart has no nonzero
zero-class object, so
$v_n=[E_n]/\norm{[E_n]}$ is defined.  After a subsequence it tends to a
unit vector $v$.  Heart positivity at the two endpoints, and at the
unperturbed middle hearts, gives in the limit
\begin{equation}\label{eq:four-halfplanes}
 \Im\bigl(e^{-2\pi i q/5}Z(v)\bigr)\ge0
 \qquad(q=0,1,2,3).
\end{equation}
These four inequalities force $Z(v)=0$.  Indeed, if $Z(v)\ne0$, choose its
argument $\vartheta\in[0,2\pi)$.  The first three inequalities force
$\vartheta\in[4\pi/5,\pi]$, whereas the fourth forces
$\vartheta\in[0,\pi/5]\cup[6\pi/5,2\pi)$, which is impossible.

Now fix any $\gamma\in\Omega$.  The middle-heart inequality at index one
holds for every $n$, so the same convergent subsequence satisfies
\[
 N_\gamma(\Phi_*^{-1}v)\ge0.
\]
Since $\gamma$ was arbitrary, this holds simultaneously for all
$\gamma\in\Omega$; there is no further subsequence or shrinking of the
middle parameter set.  Write $w=\Phi_*^{-1}v$.  The eigencharge identity
and \eqref{eq:imaginary-charge} give $Z(w)=N_0(w)=0$.  Varying
$\gamma_0,\gamma_1$ with both signs in $\Omega$ now gives
$r(w)=c(w)=0$.  The real-linear independence of the $d,e$ coefficients
of $Z$ gives $d(w)=e(w)=0$.  This contradicts $\norm{v}=1$.

Consequently some square $\Omega_1$ about zero has
$\mathcal I_{\alpha,\beta}=0$ for every pair
$\alpha,\beta\in\Omega_1$: otherwise choose counterexamples with both
parameters within distance $1/n$ of zero.  To deduce
\eqref{eq:single-shift-inclusion}, let $E\in D_3^\alpha$ and truncate in
the bounded t-structure with aisle $D_0^\beta$:
\[
 A\longrightarrow E\longrightarrow H^0_{D_0^\beta}(E)
   \longrightarrow A[1],\qquad A\in D_0^\beta[1].
\]
Here $E\in D_0^\beta$, and
$A[1]\in D_0^\beta[2]=D_5^\beta\subset D_3^\alpha$.
Extension closure therefore puts the zeroth cohomology in
$D_3^\alpha\cap(D_0^\beta[1])^\perp=0$.  Hence $E\in D_0^\beta[1]$.

Every integer is $2m+5j$.  Two such representations differ by
$(m,j)\mapsto(m+5t,j-2t)$, so $\Phi^5\simeq[2]$ proves that
\eqref{eq:fifths-aisles} is well defined.  Since
$V_1^\alpha=\Phi^3U_\alpha[-1]\subset U_\beta=V_0^\beta$,
transport by $\Phi^m[j]$ proves the first assertion of
\eqref{eq:fifths-nesting} for all $k$.  Its other assertions follow
directly from the definition.
\end{proof}

From now on $V_k=V_k^0$.  Define the heart and the block
\begin{equation}\label{eq:block-definitions}
 \mathcal H_k=V_k\cap V_{k+5}^\perp,
 \qquad \mathcal S_k=V_k\cap V_{k+1}^\perp.
\end{equation}
The first is a bounded heart because $V_{k+5}=V_k[1]$.

\begin{lemma}[Block filtrations]\label{lem:block-filtrations}
The blocks are extension closed, and
$\Hom(\mathcal S_j,\mathcal S_k)=0$ whenever $j>k$.
Every object has a finite triangle filtration with nonzero factors in
blocks of strictly decreasing indices.  In $\mathcal H_k$ there is a
torsion pair
\begin{equation}\label{eq:block-torsion-pair}
 \begin{aligned}
 \mathcal T_k=V_{k+1}\cap\mathcal H_k
     &=\langle\mathcal S_{k+4},\mathcal S_{k+3},
                \mathcal S_{k+2},\mathcal S_{k+1}\rangle,\\
 \mathcal F_k&=\mathcal S_k.
 \end{aligned}
\end{equation}
In particular, $\mathcal S_k$ is closed under subobjects in $\mathcal H_k$.
\end{lemma}

\begin{proof}
Extension closure follows from \eqref{eq:block-definitions}.  If $j>k$,
then $\mathcal S_j\subset V_j\subset V_{k+1}$ and
$\mathcal S_k\subset V_{k+1}^\perp$, proving the Hom vanishing.
Boundedness and $V_{k+5}=V_k[1]$ put any object $E$ in
$V_l\cap V_h^\perp$ for some integers $l<h$.  Truncation at $V_{l+1}$
gives
\[
 A\longrightarrow E\longrightarrow B\longrightarrow A[1],
 \quad A\in V_{l+1},\ B\in V_{l+1}^\perp.
\]
Because $E,A[1]\in V_l$, we have $B\in V_l$, hence $B\in\mathcal S_l$.
Also $B[-1]\in V_h^\perp$: for $F\in V_h$,
$F[1]\in V_h\subset V_{l+1}$ gives $\Hom(F,B[-1])=0$.
The rotated triangle thus shows $A\in V_h^\perp$.
Repeat with $A\in V_{l+1}\cap V_h^\perp$.  At index $h$ the remaining
object belongs to $V_h\cap V_h^\perp$ and is zero.  Reading the resulting
triangles from the first subobject to the last quotient gives the stated
decreasing order; zero factors are omitted.

If $E\in\mathcal H_k$, apply the same construction with $l=k$ and
$h=k+5$.  The first triangle has all its terms in $\mathcal H_k$, so it is
a short exact sequence there, with $A\in\mathcal T_k$ and
$B\in\mathcal S_k$.  The preceding Hom vanishing proves the torsion-pair
axiom.  Repeating on $A$ proves the displayed four-block description.
For completeness, a subobject $G$ of an object of $\mathcal S_k$ has
$\Hom(\mathcal T_k,G)=0$, by composing with the inclusion, so its torsion
part vanishes.  Thus $G\in\mathcal S_k$.
\end{proof}

Figure~\ref{fig:fifth-phase-grid} records the index conventions and the charge
sectors that will be established in Proposition~\ref{prop:block-support}.
\begin{figure}[tbp]
\centering
\begin{tikzpicture}[x=1.45cm,y=1cm,>=Stealth,
  every node/.style={font=\small}]
  \foreach \k in {0,...,5} {
    \node (v\k) at (\k,1.1) {$V_{\k}$};
  }
  \foreach \k in {0,...,4} {
    \pgfmathtruncatemacro{\next}{\k+1}
    \draw[->] (v\next.west) -- (v\k.east);
  }
  \draw[->] (v0.north) to[bend left=35]
    node[above] {$\Phi$} (v2.north);
  \draw[->] (v0.north) to[bend left=40]
    node[above] {$[1]$} (v5.north);
  \foreach \k in {0,...,4} {
    \fill[blue!6] (\k,-0.05) rectangle (\k+1,0.4);
    \node at (\k+0.5,0.175) {$\mathcal S_{\k}$};
  }
  \draw (0,-0.05) -- (5,-0.05);
  \foreach \k in {0,...,5} {
    \draw (\k,0.4) -- (\k,-0.12);
  }
  \node[below] at (0,-0.12) {$0$};
  \foreach \k in {1,...,4} {
    \node[below] at (\k,-0.12) {$\k/5$};
  }
  \node[below] at (5,-0.12) {$1$};
  \node[below] at (2.5,-0.65) {charge phase};
\end{tikzpicture}
\caption{The fifth-phase grid. Inclusion arrows point toward the larger
aisle. The autoequivalence $\Phi$ advances two indices, and the shift $[1]$
advances five. A nonzero object of $\mathcal S_k$ has charge phase in the
closed interval $[k/5,(k+1)/5]$, by Proposition~\ref{prop:block-support}.
At a shared endpoint, generators from both adjacent blocks enter the
definition of the slicing. The diagram shows the unperturbed grid;
Proposition~\ref{prop:fine-grid} gives the independent-parameter nesting.}
\label{fig:fifth-phase-grid}
\end{figure}


\subsection{Full numerical support on the blocks}

We first locate a block's charge using its two adjacent unperturbed hearts.
To bound its numerical class, we then place the same object in every perturbed
heart at the preceding grid position. This second step retains the two
independent inequalities that control rank and first Chern character.

\begin{proposition}\label{prop:block-support}
For the fixed Euclidean norm on $\Knum(X)_\R$, there is a constant $C>0$
such that, for every $k\in\Z$ and every nonzero $E\in\mathcal S_k$,
\begin{equation}\label{eq:block-support}
 \begin{gathered}
 Z(E)\in\{t e^{i\pi\varphi}:t>0,\ k/5\le\varphi\le(k+1)/5\},\\
 \norm{[E]}\le C|Z(E)|.
 \end{gathered}
\end{equation}
\end{proposition}

\begin{proof}
For $k=2m+5j$, the functor $\Phi^{-m}[-j]$ takes $\mathcal H_k$ to
$\mathcal A_0$ and multiplies its charge by $e^{-i\pi k/5}$.  Positivity
therefore gives $\Im(e^{-i\pi k/5}Z(E))\ge0$ on $\mathcal H_k$.
An object $E\in\mathcal S_k$ belongs both to $\mathcal H_k$ and to
\[
 \mathcal H_{k+1}[-1]
   =V_{k-4}\cap V_{k+1}^\perp.
\]
The two tests are consequently
\[
 \Im(e^{-i\pi k/5}Z(E))\ge0,\qquad
 \Im(e^{-i\pi(k+1)/5}Z(E))\le0.
\]
Their intersection is the closed sector in \eqref{eq:block-support},
together with the origin.

Choose $\varepsilon>0$ such that the closed square
$[-\varepsilon,\varepsilon]^2$ is contained in $\Omega_1$.
For every $\eta$ in this square, independent-parameter nesting gives
\[
 V_k\subset V_{k-1}^\eta,
 \qquad V_{k-1}^\eta[1]=V_{k+4}^\eta\subset V_{k+1}.
\]
Thus $E$ belongs to the heart of $V_{k-1}^\eta$ for every such $\eta$:
the second inclusion supplies exactly its required right orthogonality.
Write $k-1=2m+5j$, and set $w=[\Phi^{-m}E[-j]]$.  This is the class of an
actual object of $\mathcal A_\eta$ for every parameter under consideration,
including the shift in this formula.  Hence
\[
 N_0(w)+\eta_1c(w)+\eta_0r(w)\ge0
 \quad\text{for all }|\eta_0|,|\eta_1|\le\varepsilon.
\]
Choosing the two signs independently proves
\begin{equation}\label{eq:perturbed-support}
 \varepsilon\bigl(|r(w)|+|c(w)|\bigr)
 \le N_0(w)=\frac{\Im Z(w)}{5t_0}
 \le\frac{|Z(E)|}{5t_0}.
\end{equation}
The invertible real coordinate map
$w\mapsto(r(w),c(w),\Re Z(w),\Im Z(w))$, established in
\eqref{eq:numerical-norm-control}, now bounds $\norm{w}$ by a constant
times $|Z(E)|$.  Explicitly,
$d(w)=\Im Z(w)/(5t_0)-c(w)/2-ur(w)$, and then the real part of
\eqref{eq:central-charge} recovers $e(w)$.

We can choose $m\in\{0,1,2,3,4\}$ according to $k-1$ modulo five.
There are therefore only five norm-conversion operators $\Phi_*^m$;
the arbitrary shift $j$ only changes the sign of a numerical class.
Taking the maximum of their operator norms gives one $C$ for every $k$.
If $Z(E)=0$, the estimate gives $w=0$.  The object
$\Phi^{-m}E[-j]$ belongs to $\mathcal A_0$, so
Lemma~\ref{lem:second-tilt} forces $\Phi^{-m}E[-j]=0$ and hence $E=0$.
This excludes the origin for a nonzero block object and completes the proof.
\end{proof}

\subsection{Finite refinement inside a block}

The support estimate is now in place. We use it with the discreteness of the
full numerical lattice to obtain maximal phases and a bound on the length of
refinement; no discreteness of the charge image is needed.

Fix $k$, and put $a=k/5$, $b=(k+1)/5$.
Each nonzero $E\in\mathcal S_k$ has a unique phase
$\varphi_k(E)\in[a,b]$ with $Z(E)=|Z(E)|e^{i\pi\varphi_k(E)}$.
A subobject $F\subset E$ in $\mathcal H_k$ is called \emph{saturated in
the block} if $E/F\in\mathcal S_k$.  Its source already belongs to
$\mathcal S_k$ by Lemma~\ref{lem:block-filtrations}.  We call $E$
\emph{block-semistable} if
\[
 \varphi_k(F)\le\varphi_k(E)
 \quad\text{for every nonzero subobject saturated in the block.}
\]
The whole object is allowed as such a subobject.

\begin{lemma}[Finite block refinement]\label{lem:block-hn}
Every nonzero object of $\mathcal S_k$ has a finite filtration by short
exact sequences in $\mathcal H_k$, with all intermediate quotients in
$\mathcal S_k$, whose nonzero factors are block-semistable with strictly
decreasing phases.
\end{lemma}

\begin{proof}
Let $\ell>0$ be a lower bound for the norms of nonzero elements of the
full numerical lattice, supplied by Lemma~\ref{lem:numerical}, and define
\[
 q_k(z)=\Re\bigl(e^{-i\pi(a+b)/2}z\bigr),\qquad
 d_* =\frac{\ell\cos(\pi/10)}{C}>0.
\]
For every nonzero block object, Proposition~\ref{prop:block-support}
implies
\begin{equation}\label{eq:block-mass}
 q_k(Z(E))\ge\cos(\pi/10)|Z(E)|\ge d_*.
\end{equation}
This projection is additive in short exact sequences.  Hence every
filtration of a fixed $E$ with nonzero block factors has length at most
$q_k(Z(E))/d_*$.

If $F\subset E$ is saturated in the block, then both $F$ and $E/F$ are
block objects.  Their projections are nonnegative, so
\[
 \norm{[F]}\le C|Z(F)|
 \le \frac{C}{\cos(\pi/10)}q_k(Z(F))
 \le \frac{C}{\cos(\pi/10)}q_k(Z(E)).
\]
Only finitely many lattice classes, and therefore only finitely many
charges and phases, can occur for such $F$.  Choose a nonzero saturated
subobject $F\subset E$ of maximal phase, and among those of maximal
$|Z(F)|$.  These maxima concern finite sets of numerical values; no
finiteness of the set of subobjects is needed.

This $F$ is block-semistable.  Indeed, a saturated subobject of $F$ is
saturated in $E$, since the resulting quotient of $E$ is an extension of
two objects of $\mathcal S_k$.  If $G\subset E/F$ is a nonzero saturated
subobject, its preimage $\widetilde G\subset E$ is also saturated and is
an extension of $G$ by $F$.  Charges in a sector of width less than $\pi$
have the elementary strict see-saw property: the argument of their sum
is strictly between their arguments when those arguments differ, and
their magnitudes add when they agree.  Thus
$\varphi_k(G)>\varphi_k(F)$ contradicts maximal phase of $F$ by using
$\widetilde G$; equality contradicts maximal magnitude.  Every such $G$
therefore has phase strictly less than $\varphi_k(F)$.

Apply the same selection to the nonzero quotient $E/F$, if there is one,
and continue.  Preimages give a filtration in the original heart with all
factors and intermediate quotients in the block.  The phases strictly
decrease by the preceding argument, and \eqref{eq:block-mass} bounds the
number of nonzero factors, so the process terminates.
\end{proof}

\begin{lemma}[The unsaturated test and Hom vanishing]
\label{lem:block-hom}
If $E\in\mathcal S_k$ is block-semistable, then every nonzero subobject
$G\subset E$ in $\mathcal H_k$ satisfies
$\varphi_k(G)\le\varphi_k(E)$, without a saturation assumption.
Consequently, for block-semistable $E,F\in\mathcal S_k$,
\[
 \varphi_k(E)>\varphi_k(F)\quad\Longrightarrow\quad\Hom(E,F)=0.
\]
\end{lemma}

\begin{proof}
Suppose $G\subset E$ violates the stated inequality.  Decompose $E/G$
by \eqref{eq:block-torsion-pair} and take the preimage of its torsion
part:
\[
 0\longrightarrow T\longrightarrow E/G\longrightarrow F'\longrightarrow0,
 \qquad
 0\longrightarrow G\longrightarrow\overline G\longrightarrow T
      \longrightarrow0,
 \quad T\in\mathcal T_k,\ F'\in\mathcal S_k.
\]
The object $\overline G\subset E$ lies in $\mathcal S_k$ by subobject
closure, and is saturated because $E/\overline G=F'$.
The four-block filtration of $T$ has charge rays with phases in
$[b,a+1]$.  If $g=\varphi_k(G)$, then
\[
 a<g\le b,
 \qquad 0\le\psi-g\le a+1-g<1
 \quad\text{for every such phase }\psi.
\]
Thus $Z(G)$ and every added ray lie in the sector with phase interval
$[g,a+1]$ of width strictly less than one.  Their sum is nonzero and has
phase in that same interval.  Since $\overline G\in\mathcal S_k$, its
phase in $[a,b]$ must therefore be at least $g$.  In particular,
$\varphi_k(\overline G)>\varphi_k(E)$, a contradiction.
The strict inequality $g>a$, supplied by the alleged destabilization,
is what excludes an antipodal endpoint in this argument.  If $T=0$, the
contradiction already comes from $G$ itself being saturated.

For a nonzero map $E\to F$, take its image $I$ in $\mathcal H_k$.
Since $I\subset F$, we have $I\in\mathcal S_k$.  The kernel in $E$ is
therefore saturated.  If that kernel is zero, $\varphi_k(I)=\varphi_k(E)$;
otherwise semistability and the see-saw property give
$\varphi_k(I)\ge\varphi_k(E)$.  The unsaturated test applied to
$I\subset F$ gives $\varphi_k(I)\le\varphi_k(F)$.  These inequalities
exclude a nonzero map when the source phase is strictly larger.
\end{proof}

\subsection{The slicing and its phase cuts}

For $\varphi\in\R$, let $\mathcal P(\varphi)$ be the extension closure
of all block-semistable objects whose assigned phase equals $\varphi$.
There is one possible block unless $5\varphi\in\Z$, in which case the
two adjacent blocks are both included in this definition.

\begin{proposition}\label{prop:slicing}
The categories $\mathcal P(\varphi)$ form a slicing of $\D$ with finite
Harder--Narasimhan filtrations.  They satisfy
\begin{equation}\label{eq:slicing-covariance}
 \mathcal P(\varphi+1)=\mathcal P(\varphi)[1],\qquad
 \Phi\mathcal P(\varphi)=\mathcal P(\varphi+2/5)
 \quad(\varphi\in\R),
\end{equation}
and every nonzero $E\in\mathcal P(\varphi)$ satisfies
\begin{equation}\label{eq:semistable-support}
 Z(E)\in\R_{>0}e^{i\pi\varphi},\qquad
 \norm{[E]}\le C|Z(E)|.
\end{equation}
\end{proposition}

\begin{proof}
First consider the generators.  If two are in different blocks and the
source has strictly larger phase, its block index is larger: for $j<k$
one has $(j+1)/5\le k/5$.  Thus their Hom group vanishes by
Lemma~\ref{lem:block-filtrations}.  The case of a common block is
Lemma~\ref{lem:block-hom}.  Induction along extensions in each variable
now gives
\[
 \Hom(\mathcal P(\varphi),\mathcal P(\psi))=0
 \qquad(\varphi>\psi).
\]
The categories are additive, since finite direct sums are split
extensions.  The charge of an extension of phase-$\varphi$ generators
is a sum of positive multiples of $e^{i\pi\varphi}$ and hence is again
positive on that ray.  Likewise the triangle inequality gives
\[
 \norm{[E]}\le\sum_j\norm{[E_j]}
 \le C\sum_j|Z(E_j)|=C|Z(E)|
\]
for any such extension filtration.  This proves
\eqref{eq:semistable-support}.

Refine the decreasing block filtration of any object by
Lemma~\ref{lem:block-hn}, splicing the resulting triangles by the
octahedral axiom.  The phases are weakly decreasing globally and strictly
decreasing inside each block.  Adjacent equal phases can only occur at a
shared endpoint; merging each run of equal-phase factors into its
extension gives a finite filtration with strictly decreasing phases and
factors in the corresponding $\mathcal P(\varphi)$.  This also explains
why both blocks at a shared endpoint must enter the same phase category.

By \eqref{eq:fifths-nesting}, the functors $[1]$ and $\Phi$ take
$(\mathcal H_k,\mathcal S_k)$ to
$(\mathcal H_{k+5},\mathcal S_{k+5})$ and
$(\mathcal H_{k+2},\mathcal S_{k+2})$, respectively.  These equivalences
take exact sequences in the hearts to exact sequences, preserve the
condition of being saturated in a block, and change the assigned phase
by $1$ and $2/5$.  They therefore preserve block-semistability with these
phase increments.  Applying them and their inverses to the extension
closures proves both equalities in \eqref{eq:slicing-covariance}.
The phase-cut argument that follows proves closure under direct
summands without presupposing a slicing heart.
\end{proof}

For any interval $I\subset\R$, write $\mathcal P(I)$ for the extension
closure of the $\mathcal P(\varphi)$ with $\varphi\in I$.

\begin{lemma}[Direct construction of the phase cuts]\label{lem:phase-cuts}
For every real $c$, both pairs
\begin{equation}\label{eq:phase-cuts}
 \bigl(\mathcal P((c,\infty)),\mathcal P((-\infty,c])\bigr),
 \qquad
 \bigl(\mathcal P([c,\infty)),\mathcal P((-\infty,c))\bigr)
\end{equation}
are an aisle and its right orthogonal for a bounded t-structure.
Their hearts are respectively
$\mathcal P((c,c+1])$ and $\mathcal P([c,c+1))$.
Every interval category is the intersection of its lower and upper cut
categories, with the indicated endpoint conventions.  In particular,
these categories, including each $\mathcal P(\varphi)$, are closed under
direct summands.
\end{lemma}

\begin{proof}
Denote either pair in \eqref{eq:phase-cuts} by $(\mathcal U,\mathcal L)$.
Strict phase separation gives $\Hom(\mathcal U,\mathcal L)=0$ by
extension induction.  Equation~\eqref{eq:slicing-covariance} gives
$\mathcal U[1]\subset\mathcal U$ and
$\mathcal L[-1]\subset\mathcal L$.  Cutting the finite decreasing
filtration at $c$ and splicing triangles gives, for every $E$, a triangle
\begin{equation}\label{eq:phase-truncation}
 U_E\longrightarrow E\longrightarrow L_E\longrightarrow U_E[1],
 \qquad U_E\in\mathcal U,\ L_E\in\mathcal L.
\end{equation}
All factors of phase exactly $c$ are assigned to the same side, according
to the selected convention.  These are the t-structure axioms in the
aisle/right-orthogonal convention, and we can verify the orthogonal
characterizations directly before using any summand closure.  If
$E\in{}^\perp\mathcal L$, applying $\Hom(-,L_E)$ to
\eqref{eq:phase-truncation} shows that $\Hom(L_E,L_E)=0$, since its
neighboring terms $\Hom(U_E[1],L_E)$ and $\Hom(E,L_E)$ vanish.
Thus $L_E=0$ and $E\in\mathcal U$.  Similarly, if
$E\in\mathcal U^\perp$, apply $\Hom(U_E,-)$ and use
$\Hom(U_E,L_E[-1])=\Hom(U_E,E)=0$ to obtain $U_E=0$.
Hence
\[
 \mathcal U={}^\perp\mathcal L,\qquad
 \mathcal L=\mathcal U^\perp.
\]
They are consequently closed under direct summands, and
\eqref{eq:phase-truncation} is a genuine truncation triangle.  Finite
upper and lower bounds on the phases of each object's filtration, and
the shift law, prove boundedness.

We justify the interval description explicitly.  A nonempty decreasing
filtration with nonzero phase factors cannot have zero total object.
Indeed, the identity of its first factor maps nontrivially into the
total object: at each subsequent triangle this map remains nonzero
because $\Hom(F_1,F_j[-1])=0$.  Now group a chosen filtration of an
object at the two endpoints of an interval.  If the object lies in the
upper and lower cut categories, the preceding orthogonal
characterizations force the exterior truncation objects to be zero.
By the observation just made, their groups of factors are empty.  The
remaining factors lie in the interval.  The reverse inclusion follows
from extension closure.  This works with either choice at either
endpoint, as well as for rays and singletons.
Intersecting the aisle with the right orthogonal of its shift therefore
gives precisely the two displayed hearts.  Finally, intersections of
orthogonals are closed under summands, which proves the last assertion.
\end{proof}

\subsection{Quasi-abelian intervals and local finite length}

It remains to verify local finiteness. We first identify the exact sequences
in a short phase interval by realizing it inside a phase-cut heart. We then
use a positive projection of the charge to bound the length of every strict
filtration in that interval.

\begin{lemma}[The exact structure on a short phase interval]
\label{lem:interval-quasiabelian}
For every interval $I$ of width less than one, with any choices of open
or closed endpoints, $\mathcal P(I)$ is quasi-abelian.  Its strict exact
sequences are exactly the short exact sequences of a suitable bounded
heart whose three terms belong to $\mathcal P(I)$.
\end{lemma}

\begin{proof}
The empty interval gives the zero category.  Otherwise let $a\le b$
be its endpoints, with $b-a<1$.  If $a$ is excluded, use the heart
$\mathcal H=\mathcal P((a,a+1])$; if it is included, use
$\mathcal H=\mathcal P([a,a+1))$.  Inside this heart, the cut at $b$
gives a torsion pair $(\mathcal T,\mathcal F)$ with
$\mathcal F=\mathcal P(I)$.  Equality at $b$ is assigned to
$\mathcal F$ exactly when $b\in I$.  For example, for $I=(a,b)$ it is
\[
 \mathcal T=\mathcal P([b,a+1]),\qquad
 \mathcal F=\mathcal P((a,b))
 \quad\text{in }\mathcal P((a,a+1]).
\]
For $I=(a,b]$ the corresponding torsion class is
$\mathcal P((b,a+1])$.  The lower-closed variants use the other heart
above and the same upper-endpoint rule.  Hom orthogonality follows from
the phase inequality, and the cut triangles have all terms in
$\mathcal H$, so they give the required short exact decompositions.

We prove the quasi-abelian assertion directly for a torsion-free class
$\mathcal F$ in an abelian heart $\mathcal H$; this standard fact is also
recorded in \cite[Lemma~3.1]{Tattar2021}.
It is closed under subobjects and extensions.  For a morphism
$f:A\to B$ in $\mathcal F$, put
$K=\ker_{\mathcal H}f$, $J=\im_{\mathcal H}f$, and
$Q=\operatorname{coker}_{\mathcal H}f$.
The kernel in $\mathcal F$ is $K$, and the cokernel is
$Q/t(Q)$, where $t(Q)$ denotes the torsion part of $Q$.
Indeed, maps from $Q$ to an object of $\mathcal F$ annihilate $t(Q)$;
this proves the cokernel's universal property.  Both $K$ and $J$ lie in
$\mathcal F$.  Consequently the internal coimage of $f$ is $J$, whereas
its internal image is the preimage of $t(Q)$ under $B\to Q$.

Recall that a morphism is strict if its canonical coimage-to-image map
is an isomorphism.  The preceding formulas show that a strict
monomorphism in $\mathcal F$ is exactly a monomorphism in $\mathcal H$
whose cokernel belongs to $\mathcal F$.  A strict epimorphism in
$\mathcal F$ is exactly an epimorphism in $\mathcal H$ between its
objects: for an internal epimorphism the internal image is $B$, and
strictness says precisely that $J=B$.  This statement is about strict
epimorphisms, not all epimorphisms of the subcategory.

The pushout in $\mathcal H$ of a strict monomorphism $A\hookrightarrow B$
along $A\to A'$ fits into
\[
 0\longrightarrow A'\longrightarrow B'\longrightarrow B/A
     \longrightarrow0.
\]
Both outside terms are in $\mathcal F$, so $B'\in\mathcal F$ and the
new monomorphism is strict.  The pullback in $\mathcal H$ of a strict
epimorphism $A\twoheadrightarrow B$ along $B'\to B$ is a subobject of
$A\oplus B'$, hence lies in $\mathcal F$, and its projection to $B'$ is
an ambient epimorphism, hence strict.  These ambient universal
properties remain universal in the full subcategory $\mathcal F$.
Thus kernels and cokernels exist, pushouts preserve strict
monomorphisms, and pullbacks preserve strict epimorphisms.  This proves
quasi-abelianity and the stated description of strict exact sequences.
\end{proof}

\begin{proposition}[Local finiteness]\label{prop:local-finiteness}
Every phase interval of width less than one generates a finite-length
quasi-abelian category.  Hence the slicing in
Proposition~\ref{prop:slicing} is locally finite.
\end{proposition}

\begin{proof}
Let $I$ have endpoints $a\le b$ and width $w=b-a<1$, and put
\[
 q_I(z)=\Re\bigl(e^{-i\pi(a+b)/2}z\bigr),\qquad
 d_I=\frac{\ell}{C}\cos(\pi w/2)>0,
\]
where $\ell$ is the full-lattice norm lower bound used in
Lemma~\ref{lem:block-hn}.  A nonzero semistable $F$ of phase in $I$ has
nonzero numerical class, since its charge is nonzero.  Therefore
\eqref{eq:semistable-support} gives
\[
 q_I(Z(F))\ge\cos(\pi w/2)|Z(F)|\ge d_I.
\]
Every nonzero object of $\mathcal P(I)$ is a finite extension of nonzero
semistable objects with phases in $I$.  Additivity thus gives the same
lower bound $q_I(Z(E))\ge d_I$ for every such object.

By Lemma~\ref{lem:interval-quasiabelian}, a strict exact sequence in
$\mathcal P(I)$ is exact in its ambient heart, so the projection $q_I Z$
is additive on it.  Every nonzero successive quotient in a strict
filtration of a fixed $E\in\mathcal P(I)$ is again an object of
$\mathcal P(I)$ and consumes at least $d_I$ of this projection.
The number of such quotients is consequently at most
$q_I(Z(E))/d_I$.  This bounds both increasing and decreasing strict
subobject chains, and proves finite length in the quasi-abelian exact
structure.  The argument uses discreteness of the full numerical
lattice; it requires no discreteness of its image under $Z$.
\end{proof}

The charge is numerical and normalized on point sheaves by
Proposition~\ref{prop:central-charge}; its eigencharge identity combines
with \eqref{eq:slicing-covariance}.  Proposition~\ref{prop:slicing},
Lemma~\ref{lem:phase-cuts}, and Proposition~\ref{prop:local-finiteness}
give all slicing, Harder--Narasimhan, and local-finiteness requirements.
Finally \eqref{eq:semistable-support} is the support property on the
\emph{full} space $\Knum(X)_\R$.  This completes the proof of
Theorem~\ref{thm:main}.
